% ====================================================================
% Chapter 2: Theoretical Foundations and Taxonomy
% ====================================================================

\section{Theoretical Foundations and Taxonomy of Self-Evolving Agents}\label{ch:theory}

\subsection{Chapter Overview}\label{sec:ch2-overview}

The theoretical foundations of Agent Self-Evolution draw from four converging research traditions. Evolutionary computation provides the ``generate--vary--select--retain'' framework for open-ended search over non-differentiable spaces. Self-reference and G\"{o}del machines offer a recursive lens through which a system can modify its own structure and empirically verify improvements. Meta-learning and self-training supply the machine-learning formalism for ``learning how to learn'' across task distributions. Reinforcement learning, online learning, and program synthesis contribute mathematical characterizations of objectives, policies, environments, evaluators, and update operators. This chapter does not claim that a mature unified theory already exists---the literature is still coalescing rapidly. Instead, the goal is to abstract the common structures underlying diverse works, providing conceptual coordinates for the method taxonomy in Chapter~\ref{ch:methods} and the system analysis in subsequent chapters.

A central contribution of this survey is the identification of five recurring ``evolution loops'' that appear, in various combinations, across the systems we study. These loops are not mutually exclusive; most real systems instantiate two or more simultaneously. However, isolating them as analytical primitives clarifies design choices, trade-offs, and failure modes. The five loops are: (1) the \emph{Specification-to-Execution Loop}, which converts natural-language goals into runnable pipelines; (2) the \emph{Search Loop}, which explores architectures, prompts, code, agent designs, and hyperparameters; (3) the \emph{Evaluator Loop}, which scores candidates through tests, benchmarks, verifiers, or deployment metrics; (4) the \emph{Reflection Loop}, which transforms failures into memories, feedback signals, or revised candidate generators; and (5) the \emph{Population Loop}, which maintains multiple candidates and applies selection, mutation, recombination, and specialization across them. Sections~\ref{sec:theory-evo} through~\ref{sec:theory-formal} develop the theoretical machinery; Section~\ref{sec:theory-five-loops} provides the detailed loop taxonomy.

\subsection{Evolutionary Computation and LLM Integration}\label{sec:theory-evo}

\subsubsection{Evolutionary Computation as Search Framework}

The foundational idea of evolutionary computation (EC) is straightforward: maintain a population or archive of candidate solutions in a search space, and iteratively improve them through variation, recombination, and selection based on a fitness function. Classical genetic algorithms, genetic programming, evolution strategies, novelty search, and quality-diversity algorithms all belong to this tradition. Their strength lies in requiring no differentiability of the objective function and no fixed-length representation of candidates---as long as one can \emph{evaluate} a candidate, search can proceed.

Agent Self-Evolution operates in search spaces that are precisely the kind EC handles well. Prompts are discrete text strings; tool chains are combinatorial structures; agent code is executable programs; memories are natural-language or vector collections; multi-agent topologies are graphs; world models and evaluators may themselves be composed of code or LLM prompts. None of these objects admits gradient-based optimization in the usual sense. Evolutionary computation thus provides a natural outer-loop optimization framework for self-evolving agents.

\subsubsection{LLM as Semantic Mutator}

Large language models transform two critical components of the evolutionary loop. First, the LLM acts as a \emph{strong semantic mutator}. Traditional variation operators---bit flips, Gaussian perturbations, subtree crossovers---are random or near-random. An LLM, by contrast, can generate \emph{targeted} modifications based on failure explanations: adding a test-running step for a coding agent, rewriting the exploration strategy for a web agent, inserting a self-verification sub-routine for a reasoning agent, or designing a debate protocol for a multi-agent system. ADAS's Meta Agent Search, DGM's self-modification, SICA's code updates, and AlphaEvolve's algorithm candidate generation all use the LLM as a high-level \emph{program mutation operator}, producing changes informed by context rather than blind perturbation.

Second, the LLM can participate in \emph{evaluation and explanation}. Self-Rewarding Language Models, Meta-Rewarding, IterAlign, ACE, and EvolveR all employ the model to generate scores, principles, critiques, or policy summaries. This means selection pressure need not come solely from hard-coded metrics; it can also derive from natural-language preferences, process judgments, and qualitative feedback. However, using LLMs as judges introduces well-documented biases: length preference, position bias, verbosity rewards, and---most critically---the risk that the evaluator and the generator share the same failure modes, producing self-consistent but systematically wrong assessments.

\subsubsection{The Archive as Externalized Search Experience}

The interaction between LLMs and evolutionary computation also introduces new design challenges. LLM-generated candidates carry strong priors from training data, so the search is not a uniform exploration of the space. The archive---the collection of previously evaluated candidates---plays a critical role. If it retains only the highest-scoring candidates, the system may converge prematurely; if it retains everything, context length and computational cost explode. DGM makes the deliberate choice to maintain a growing archive of agent variants, embracing open-ended exploration to avoid local optima. ADAS leverages historical discoveries within the meta-agent's context to compose new agent designs. Both illustrate a key principle: in Agent Self-Evolution, the archive is not merely a log but the \emph{externalized carrier of reusable search experience}.

\subsubsection{Five Design Dimensions}

From the evolutionary computation perspective, Agent Self-Evolution involves at least five design dimensions:
\begin{enumerate}
\item \textbf{Representation}: What constitutes a candidate individual? Options include prompts, policies, memory states, tool plans, Python code, workflow graphs, or model weights.
\item \textbf{Variation}: How are new candidates generated? Mechanisms range from random perturbation, LLM rewriting, reflection-based revision, textual gradients, MCTS-guided search, RL updates, to program patches.
\item \textbf{Fitness}: How are candidates scored? Signals include benchmark performance, unit tests, environment rewards, judge scores, human preferences, cost constraints, or safety rule compliance.
\item \textbf{Selection}: Which candidates survive? Strategies include keeping the best, keeping the novel, quality-diversity archiving, task-domain stratification, or rollback gating.
\item \textbf{Inheritance}: How is experience transferred across generations? Carriers include model weights, in-context examples, memory banks, code repositories, skill libraries, principle collections, or archive metadata.
\end{enumerate}

\subsection{Self-Reference and G\"{o}del Machines}\label{sec:theory-godel}

\subsubsection{Theoretical Background}

Self-reference is the most theoretically provocative aspect of Agent Self-Evolution. The concept traces to Schmidhuber's G\"{o}del machine: a theoretical system containing a problem solver and a proof searcher. When the proof searcher finds a self-modification that provably improves expected utility without violating global optimality, the system executes the modification. The importance of this framework is conceptual: it transforms ``recursive self-improvement'' from science fiction into a formal problem---how does a system represent itself, reason about the consequences of its own modifications, and avoid local improvements that compromise long-term objectives?

\subsubsection{From Formal Proofs to Empirical Verification}

Modern LLM-based agents deviate significantly from the original G\"{o}del machine. G\"{o}del Agent, DGM, and SICA all borrow the self-reference idea, but they replace formal proofs with empirical evaluation. G\"{o}del Agent allows the LLM to modify its own logic through runtime monkey-patching. SICA lets a coding agent autonomously edit its own source code to improve performance on SWE-Bench Verified. DGM maintains an executable archive of agent variants, where stronger coding agents generate even stronger descendants. These systems demonstrate the possibility of ``practical G\"{o}delization'': rather than requiring the system to prove global optimality, it suffices that each variant is \emph{reliably} better in isolated evaluation.

\subsubsection{The Verification Gap}

This shift from formal proof to empirical evaluation raises a core theoretical question: when formal proofs are unavailable, what level of empirical verification suffices to support self-modification? If the evaluator's coverage is insufficient, the system may optimize for gaps in the evaluation rather than genuine improvement. If the evaluation tasks are not representative of real-world objectives, the system may Goodhart the proxy. If the agent can modify the evaluator or bypass tests, the ``improvement'' is meaningless.

Consequently, empirical G\"{o}del agents must decompose self-reference into controlled layers: a \emph{modifiable layer} (prompts, tools, code, strategies), an \emph{immutable evaluation layer} (isolated evaluators that the agent cannot directly modify), a \emph{sandbox execution layer} (isolated environments for testing), an \emph{archive recording layer} (complete lineage and rollback paths), and a \emph{human audit layer} (for reviewing high-risk modifications). Without these layers, self-reference easily degenerates into uncontrolled self-rewriting.

\subsubsection{Positive Feedback Between Capability and Self-Modification}

Self-reference also involves a recursive ``capability feedback'' structure. DGM's key insight is that stronger coding agents can not only solve more external tasks but also better improve their own code. This creates positive feedback between downstream task performance and self-modification capability---a structure stronger than ordinary AutoML, because the searcher itself is the search target. ADAS retains a fixed meta-agent, while DGM allows the improved agent to assume self-improvement responsibility. This distinction corresponds to two theoretical models: the \emph{exogenous optimizer model} (easier to control) and the \emph{endogenous optimizer model} (closer to recursive self-improvement, but with higher safety and evaluation difficulty).

\subsection{Meta-Learning and the Theory of Self-Improvement}\label{sec:theory-meta}

\subsubsection{Learning to Learn Across Tasks}

Meta-learning concerns ``learning how to learn.'' In Agent Self-Evolution, meta-learning need not take the form of MAML-style gradient updates; it can also manifest as the progressive optimization of experience principles, in-context strategies, memory retrieval patterns, task curricula, and tool selections. Reflexion compresses failure trajectories into natural-language reflections---effectively storing a first-order policy update in context. ExpeL integrates experiences across multiple tasks into transferable insights. ICE (Investigate-Consolidate-Exploit) converts task trajectories into simplified workflows. ACE treats the context as an evolving playbook. EvolveR organizes offline self-distillation, online interaction, and policy evolution into a lifecycle.

These works extend ``learning'' from the weight space to the space of external cognitive artifacts---prompts, memories, skill libraries, principle collections, and world models.

\subsubsection{Three-Layer Policy Formalization}

An LLM-based agent can be represented as a three-layer policy:
\begin{itemize}
\item \textbf{Base layer} $\theta$: The foundation model parameters, obtained through pre-training and post-training alignment.
\item \textbf{Behavioral interface} $\varphi$: The editable behavioral layer, including system prompts, tool descriptions, planners, critics, code modules, and multi-agent topologies.
\item \textbf{Experience state} $m$: The accumulated experience, including memories, skill libraries, principle collections, failure reflections, world models, and archive metadata.
\end{itemize}

Traditional fine-tuning primarily updates $\theta$; prompt engineering updates $\varphi$; memory-based self-evolution updates $m$; systems like SICA, DGM, and ADAS simultaneously update $\varphi$ and code-level structures; RAGEN, Absolute Zero, and SPIRAL may update $\theta$ or policy parameters. The generalized policy can be written as $\pi(a \mid s; \theta, \varphi, m)$, and self-evolution is the process of updating components of $(\theta, \varphi, m)$ under the drive of feedback $r$ or evaluation $e$.

\subsubsection{Task Distribution and Transfer}

The meta-learning framework emphasizes task distributions. A system that iterates on a single fixed task may only be performing local optimization; extracting cross-task strategies is necessary for genuine self-evolution. Voyager's skill library is reusable across Minecraft exploration tasks. ExpeL demonstrates forward transfer of extracted insights. ADAS reports that discovered agent designs transfer across domains and models. ReasoningBank distills reasoning strategies from both successful and failed episodes for use on new problems. Theoretically, this can be expressed as minimizing expected loss over a task distribution $\mathcal{T}$: the system does not merely seek the optimal behavior for a single task $\tau$ but learns an update operator $U$ such that, after a few interactions, it performs better on new tasks.

\subsubsection{Online Learning View}

Self-improvement can also be viewed as online learning. At each round $t$, the agent executes a trajectory $\tau_t$ in environment $\mathcal{E}$, the evaluator provides feedback $f_t$, and the updater $U$ produces a new state $z_{t+1} = U(z_t, \tau_t, f_t)$, where $z_t = (\theta_t, \varphi_t, m_t)$. If $\theta$ is fixed and $m$ is updated, this is memory-based self-evolution; if $\varphi$ is updated, it is prompt/architecture self-evolution; if $\theta$ is updated, it is training-based self-evolution; if $\mathcal{E}$ or the task generator is also updated, it is self-curriculum or self-play. This unified framework facilitates comparison: Reflexion's $U$ writes to reflective memory, Self-Rewarding's $U$ executes Iterative DPO, RAGEN's $U$ is trajectory-level RL, DGM's $U$ is code patch with archive selection, and ACE's $U$ is incremental context playbook restructuring.

% ====================================================================
% Figure 2.1: Unified Formalized Evolution Loop (TikZ)
% ====================================================================
\begin{figure}[htbp]
\centering
\begin{tikzpicture}[
  state/.style={rectangle, rounded corners=6pt, draw=#1!70!black, fill=#1!12,
    text width=3.2cm, minimum height=1.3cm, align=center, font=\small\bfseries},
  op/.style={rectangle, rounded corners=3pt, draw=gray!60, fill=gray!6,
    text width=2.6cm, minimum height=0.9cm, align=center, font=\scriptsize},
  arr/.style={-{Stealth[length=2.5mm]}, thick, color=gray!70},
  darr/.style={-{Stealth[length=2mm]}, thick, dashed, color=gray!50},
  mathlbl/.style={font=\scriptsize, color=gray!60!black, align=center, fill=white,
    inner sep=1pt},
]
  % State vector box
  \node[rectangle, rounded corners=8pt, draw=cat1!50!black, fill=cat1!5,
    text width=5cm, minimum height=1cm, align=center, font=\small]
    (state) at (0, 1.8)
    {System State $z_t = (\theta,\, c,\, g,\, m,\, \mathcal{A})$};

  % Main loop nodes
  \node[state=cat1] (gen) at (-5, -0.8) {Generate};
  \node[state=cat2] (eval) at (0, -0.8) {Evaluate};
  \node[state=cat3] (sel) at (5, -0.8) {Select/Retain};
  \node[state=cat4] (upd) at (5, -3.8) {Vary/Update\\$U_k$};
  \node[state=cat5] (arc) at (0, -3.8) {Archive\\$\mathcal{A}_{t+1}$};
  \node[state=cat6] (safe) at (-5, -3.8) {Safety Gate\\$C(x_t, z_t)$};

  % Arrows forming the loop
  \draw[arr] (gen) -- (eval)
    node[midway, above, mathlbl] {Trajectory $x_t \sim \pi_z(\cdot|\tau_t, \mathcal{E})$};
  \draw[arr] (eval) -- (sel)
    node[midway, above, mathlbl] {Feedback $y_t = V(x_t, \tau_t, \mathcal{E})$};
  \draw[arr] (sel) -- (upd)
    node[midway, right, mathlbl] {Candidate set\\$\{U_k(z_t, x_t, y_t)\}_{k=1}^K$};
  \draw[arr] (upd) -- (arc)
    node[midway, below, mathlbl] {Diversity $D$};
  \draw[arr] (arc) -- (safe)
    node[midway, below, mathlbl] {Gating};
  \draw[arr] (safe) -- (gen)
    node[midway, left, mathlbl] {Update $z_{t+1}$};

  % State to loop connections
  \draw[darr, cat1!50] (state.south) -- ++(0, -0.3) -| (gen.north);
  \draw[darr, cat2!50] (state.south) -- (eval.north);

  % External inputs
  \node[op] (task) at (-5, 1.8) {Task distribution\\$P(\tau)$, Env $\mathcal{E}$};
  \draw[darr, cat1!40] (task.east) -- (state.west);

  % Method examples along the bottom
  \node[rectangle, rounded corners=3pt, draw=cat2!50, fill=cat2!5,
    text width=10cm, align=center, font=\scriptsize, inner sep=3pt]
    (examples) at (0, -5.6)
    {Reflexion: $U$ writes reflection $\to m$ \quad
     Self-Rewarding: $U$ = DPO $\to \theta$ \quad
     DGM: $U$ = code patch $\to g,\mathcal{A}$ \quad
     ACE: $U$ = playbook restructure $\to c$};

  \draw[darr, gray!40] (safe.south) -- ++(0,-0.5) -- (examples.west);
  \draw[darr, gray!40] (upd.south) -- ++(0,-0.5) -- (examples.east);

\end{tikzpicture}
\caption{Unified formalized evolution loop for Agent Self-Evolution. The system state $z_t=(\theta,c,g,m,\mathcal{A})$ is iteratively updated through a ``generate--evaluate--select--vary'' cycle, jointly modulated by task distribution $P(\tau)$, evaluator $V$, and safety constraints $C$. The bottom row maps specific methods to their $U$ operator.}
\label{fig:ch2-evolution-loop}
\end{figure}

\subsection{Mathematical Formalization of Self-Evolution}\label{sec:theory-formal}

\subsubsection{Unified State and Update Equations}

For unified description, this survey adopts the following formalization. Let the agent system state be $z = (\theta, c, g, m, \mathcal{A})$, where $\theta$ represents the base model or trainable parameters, $c$ represents the context and prompts, $g$ represents tools, code, and control-flow graphs, $m$ represents memories, skills, principles, and world models, and $\mathcal{A}$ represents the archive or collection of historical candidates. Given a task distribution $P(\tau)$, environment $\mathcal{E}$, evaluator $V$, and safety constraints $C$, the agent generates a behavioral trajectory $x_t \sim \pi_z(\cdot \mid \tau_t, \mathcal{E})$ at round $t$, receives feedback $y_t = V(x_t, \tau_t, \mathcal{E})$, and undergoes constraint checking $C(x_t, z_t)$. The updater $U$ produces a set of candidate states $\{z'_{t,k}\}$, and the selector $S$ chooses $z_{t+1}$ or adds multiple candidates to $\mathcal{A}$ based on quality, diversity, cost, and safety gating:
\begin{equation}\label{eq:update}
z_{t+1},\, \mathcal{A}_{t+1} = S\!\Big(\mathcal{A}_t \cup \{U_k(z_t, x_t, y_t)\}_{k=1}^K ;\; V,\, C,\, D\Big)
\end{equation}
where $D$ represents a diversity or novelty measure. Different methods primarily differ in the definitions of $U$, $V$, $S$, and $D$.

\subsubsection{Mapping Methods to Formal Components}

Self-Refine's $U$ rewrites the output using self-feedback from the same model, with state changes primarily occurring on a single sample's text. Reflexion's $U$ writes reflections to $m$. Self-Rewarding's $V$ is the model's own judge, and $U$ is DPO. RAGEN's $V$ comes from multi-turn environment rewards, and $U$ is StarPO. ADAS's $U$ has the meta-agent write new agent code, changing $g$. DGM's $U$ has the current agent modify its own code, and $\mathcal{A}$ preserves lineage. AlphaEvolve's $V$ is a programmatic evaluator, and $U$ generates algorithm code. Memory-R1's $U$ learns when to write, update, delete, and retrieve memories.

\subsubsection{Hierarchy of Improvement}

Formalization must also distinguish levels of ``improvement.'' \emph{Local improvement} means $V$ increases on the current task set. \emph{Robust improvement} requires increase on an independent test set. \emph{Transfer improvement} requires increase on a new task distribution or new model. \emph{Open-ended improvement} requires the archive to continuously produce novel and useful stepping stones. \emph{Safe improvement} requires capability growth without increasing constraint violations. Many papers demonstrate only the first two; works like ADAS, DGM, Voyager, and ReasoningBank attempt to demonstrate transfer or open-endedness. Future theory needs to more rigorously distinguish these levels; otherwise, ``self-improvement'' is easily over-interpreted from a single benchmark improvement.

\subsubsection{Selection Pressure and Multi-Objective Optimization}

Another key concept is \emph{selection pressure}. The evaluator $V$ determines where the system evolves. If $V$ is unit tests, the system optimizes for test passage. If $V$ is an LLM judge, the system optimizes for judge preferences. If $V$ is environment reward, the system exploits reward structure loopholes. If $V$ is human preference, the system may overfit the annotator distribution. Goodhart's law applies: when a metric becomes a target, it ceases to be a good measure. Self-evolving systems therefore typically require multi-objective optimization: maximize task reward $R$, minimize cost $\mathrm{Cost}$, satisfy safety constraints $C$, maintain diversity $D$, and control complexity $\Omega$. This can be expressed as finding $z$ that maximizes:
\begin{equation}\label{eq:multi-obj}
\mathbb{E}[R(z)] - \lambda\,\mathrm{Cost}(z) + \beta\,D(z) - \gamma\,\Omega(z)
\end{equation}
subject to $C(z) = \mathrm{true}$. Practical systems approximate this through gating, ablation, rollback, and human review.

\subsubsection{Theoretical Boundaries}

The theoretical boundaries of self-evolution must also be acknowledged. First, without external or independent evaluators, self-rewarding may collapse into self-consistency illusion. Second, if the task distribution is too narrow, the system may learn benchmark-specific hacks. Third, if the modifiable scope is too broad and constraints are insufficient, the system may corrupt evaluation, logging, or safety boundaries. Fourth, if the archive lacks curation, long-term memory accumulates noise, contradictions, and context collapse. Fifth, if the underlying foundation model lacks sufficient capability, LLM-generated variations may be unable to cross critical technical thresholds. Agent Self-Evolution is not a magic growth engine; it is an optimization process that couples base model capabilities, external feedback, search mechanisms, and engineering constraints. It amplifies existing strengths but also amplifies deficiencies in the evaluator and environment design.

% ====================================================================
% Figure 2.2: Classical Evolution vs. Agent Evolution Comparison
% ====================================================================
\begin{figure}[htbp]
\centering
\begin{tikzpicture}[
  header/.style={rectangle, rounded corners=6pt, draw=#1!70!black, fill=#1!20,
    text width=11cm, minimum height=1cm, align=center, font=\bfseries},
  rowlabel/.style={rectangle, rounded corners=3pt, draw=gray!50, fill=gray!8,
    text width=2.8cm, minimum height=1cm, align=center, font=\small\bfseries},
  cell/.style={rectangle, rounded corners=2pt, draw=gray!30, fill=#1,
    text width=4.5cm, minimum height=1cm, align=center, font=\scriptsize},
  divider/.style={draw=gray!40, dashed, thick},
]
  % Headers
  \node[header=gray] (title) at (0, 0) {Classical Evolutionary Computation vs.\ Agent Self-Evolution};

  % Column headers
  \node[rowlabel, fill=cat4!15, draw=cat4!50] (col1) at (-4.5, -1.5) {Dimension};
  \node[rowlabel, fill=cat1!15, draw=cat1!50] (col2) at (0, -1.5) {Classical EC};
  \node[rowlabel, fill=cat5!15, draw=cat5!50] (col3) at (4.5, -1.5) {Agent Self-Evolution};

  % Row 1: Representation
  \node[rowlabel] (r1) at (-4.5, -2.8) {Representation};
  \node[cell=cat1!5] (c1) at (0, -2.8) {Fixed-length gene strings\\Real vectors / S-expressions};
  \node[cell=cat5!5] (d1) at (4.5, -2.8) {Prompts, policies, memories,\\Python code, workflow graphs, $\theta$};

  % Row 2: Mutation
  \node[rowlabel] (r2) at (-4.5, -3.9) {Mutation};
  \node[cell=cat1!5] (c2) at (0, -3.9) {Random flips, crossover,\\Gaussian perturbation};
  \node[cell=cat5!5] (d2) at (4.5, -3.9) {LLM rewrite, reflection text,\\textual gradient, RL update, patch};

  % Row 3: Fitness
  \node[rowlabel] (r3) at (-4.5, -5.0) {Fitness};
  \node[cell=cat1!5] (c3) at (0, -5.0) {Predefined scalar function\\Hand-crafted};
  \node[cell=cat5!5] (d3) at (4.5, -5.0) {Benchmarks, unit tests,\\LLM judges, human preferences};

  % Row 4: Selection
  \node[rowlabel] (r4) at (-4.5, -6.1) {Selection};
  \node[cell=cat1!5] (c4) at (0, -6.1) {Tournament / roulette\\Fixed strategy};
  \node[cell=cat5!5] (d4) at (4.5, -6.1) {Quality-diversity, rollback gating,\\safety constraints $C$, cost control};

  % Row 5: Self-reference
  \node[rowlabel] (r5) at (-4.5, -7.2) {Self-reference};
  \node[cell=cat1!5] (c5) at (0, -7.2) {None\\Searcher and searchee separated};
  \node[cell=cat5!5] (d5) at (4.5, -7.2) {G\"{o}delization / empirical self-mod.\\Searcher itself is the search target};

  % Row 6: Theory
  \node[rowlabel] (r6) at (-4.5, -8.3) {Theory};
  \node[cell=cat1!5] (c6) at (0, -8.3) {Convergence theorems\\No-Free-Lunch};
  \node[cell=cat5!5] (d6) at (4.5, -8.3) {Primarily empirical\\Formal framework emerging};

  % Horizontal divider between header and rows
  \draw[divider] (-7.2, -2.1) -- (7.2, -2.1);
  % Vertical dividers
  \draw[divider] (-2.2, -1.5) -- (-2.2, -8.8);
  \draw[divider] (2.2, -1.5) -- (2.2, -8.8);

  % Bottom annotation
  \node[font=\scriptsize, text=gray!50!black, align=center] at (0, -9.2)
    {Core differences: LLM as semantic mutator and judge, the searcher is itself searchable, improvements must be auditable.};

\end{tikzpicture}
\caption{Systematic comparison between classical evolutionary computation and Agent Self-Evolution across six dimensions: representation, mutation, fitness, selection, self-reference, and theoretical guarantees.}
\label{fig:ch2-comparison}
\end{figure}


% ====================================================================
% Table 2.1: Theory Components to Engineering Mechanisms Mapping
% ====================================================================
\begin{table}[htbp]
\centering
\small
\caption{Mapping between Agent Self-Evolution theory components and engineering mechanisms.}
\label{tab:ch2-theory-components}
\begin{tabular}{p{0.22\textwidth}p{0.30\textwidth}p{0.36\textwidth}}
\toprule
Theory Component & Engineering Carrier & Key Risk \\
\midrule
Evolutionary search & Candidate prompts, workflows, agent code, and archive & Premature convergence, evaluation cost explosion, diversity loss \\
Self-reference & Self-modification patches, runtime hooks, versioned agent variants & Local improvements breaking long-term objectives, privilege escalation, audit difficulty \\
Meta-learning & Experience libraries, principle collections, skill libraries, cross-task initialization & Negative transfer, overfitting to historical tasks, context pollution \\
RL / Online learning & Rewards, judges, unit tests, environment feedback, and rollback gating & Reward hacking, judge bias, non-stationary feedback \\
Software engineering constraints & Sandboxes, CI, regression testing, logging, and safety policies & Insufficient test coverage, hidden regressions, irreproducibility \\
\bottomrule
\end{tabular}
\end{table}

\subsection{Summary}\label{sec:theory-summary}

The theoretical landscape of Agent Self-Evolution can be summarized as follows. LLMs provide semantic generation and reasoning capabilities. Evolutionary computation provides open-ended candidate search. Self-reference provides the recursive self-improvement objective. Meta-learning provides cross-task experience reuse. Reinforcement learning and online learning provide feedback-driven updates. Software engineering provides testing, sandboxing, version control, and rollback mechanisms.

The field does not yet possess a unified theory as mature as that of supervised learning, but it has established a stable set of questions: \emph{What is being changed?} (the five state components $\theta, c, g, m, \mathcal{A}$). \emph{Is the evaluator reliable?} (the Evaluator Loop and its failure modes). \emph{How is the archive managed?} (the Population Loop and stepping stones). \emph{Can improvements transfer?} (the meta-learning dimension). \emph{Are safety constraints tamper-proof?} (the G\"{o}del machine verification gap).

The five evolution loops---Specification-to-Execution, Search, Evaluator, Reflection, and Population---provide the taxonomic framework for answering these questions systematically. Chapter~\ref{ch:methods} will map specific methods to these loops, and subsequent chapters will analyze how real systems compose them.

\subsection{The Five Evolution Loops: A Unified Taxonomy}\label{sec:theory-five-loops}

Building on the evolutionary computation framework above, we now present the five evolution loops that constitute our central taxonomic contribution. Each loop represents a recurring pattern observed across multiple systems, and each loop has characteristic inputs, outputs, internal state, and failure modes. Real systems typically compose multiple loops; the taxonomy is analytical, not prescriptive.

\subsubsection{The Specification-to-Execution Loop}\label{sec:loop-spec}

\textbf{Definition.} The Specification-to-Execution Loop converts a natural-language or high-level specification of a goal into a runnable, evaluated artifact---a machine learning pipeline, an agent workflow, an executable program, or a deployed service. The loop iterates between specification interpretation, candidate generation, execution, and specification refinement.

\textbf{Representative Systems.} AutoML-GPT uses GPT as a natural-language bridge to convert task cards, model cards, and data cards into data processing, model architecture, hyperparameter optimization, and training log prediction. AutoML-Agent decomposes full-pipeline AutoML into specialized agents for data retrieval, preprocessing, modeling, evaluation, and deployment, combining retrieval-augmented planning with multi-stage verification. The Text-to-ML paradigm treats ML pipeline generation as autonomous program synthesis combined with search. ADAS takes this further: rather than generating a single pipeline, the meta-agent searches over the \emph{space of agent designs}, producing novel architectures that are then evaluated on held-out tasks.

\textbf{Formal Structure.} Let the specification be a natural-language description $s \in \mathcal{S}$. The generator $G$ produces a candidate artifact $a = G(s; h)$, where $h$ represents historical context (previous attempts, similar tasks). The executor $E$ runs $a$ in an environment $\mathcal{E}$, producing a trace $\tau = E(a, \mathcal{E})$. The evaluator $V$ scores the trace: $v = V(\tau, s)$. The specification refiner $R$ updates the specification or the context: $s' = R(s, \tau, v)$. The loop continues until $v$ exceeds a threshold or a budget is exhausted:
\begin{equation}\label{eq:spec-loop}
s_{t+1} = R\big(s_t,\; E(G(s_t; h_t), \mathcal{E}),\; V(\tau_t, s_t)\big)
\end{equation}

\textbf{Key Challenges.} Ambiguity in natural-language specifications is the primary obstacle. The same specification can yield wildly different artifacts depending on interpretation. Current systems address this through structured task cards (AutoML-GPT), multi-agent debate over interpretation (AutoML-Agent), or iterative clarification. However, the gap between what a user \emph{means} and what the system \emph{understands} remains the weakest link. Additionally, the cost of executing full pipelines on each iteration can be prohibitive, motivating approximate or partial evaluation strategies.

\textbf{Connection to Other Loops.} The Specification-to-Execution Loop naturally chains into the Evaluator Loop (Section~\ref{sec:loop-evaluator}), since execution produces traces that must be scored. It also connects to the Reflection Loop (Section~\ref{sec:loop-reflection}), since specification refinement is itself a form of reflective improvement.

\subsubsection{The Search Loop}\label{sec:loop-search}

\textbf{Definition.} The Search Loop explores a space of candidates---prompts, architectures, code, agent designs, or hyperparameters---guided by feedback from evaluation. Unlike the Specification-to-Execution Loop, which focuses on converting goals to artifacts, the Search Loop focuses on \emph{systematically exploring and improving} the space of possible solutions.

\textbf{Representative Systems.} OPRO treats the LLM itself as an optimizer: the prompt describes the optimization task, and the model generates new candidate solutions conditioned on the history of past solutions and their scores. AlphaEvolve combines Gemini Flash for fast candidate generation with Gemini Pro for deep evaluation, using MAP-Elites quality-diversity search to maintain a diverse archive of algorithm candidates. OpenEvolve extends FunSearch's program evolution framework with more sophisticated variation operators. EvoPrompting uses language models as evolutionary operators over architecture code, evaluating candidates on algorithmic reasoning tasks. ADAS searches over the space of agent architectures using a meta-agent that proposes and evaluates new designs.

\textbf{Formal Structure.} The Search Loop maintains a population or archive $\mathcal{A}_t$ of candidates with associated fitness scores. At each iteration, a selection mechanism $\sigma$ chooses parent candidates; a variation operator $\mu$ (powered by the LLM) generates offspring; an evaluator $V$ scores the offspring; and an update rule $\phi$ determines which candidates are retained:
\begin{equation}\label{eq:search-loop}
\mathcal{A}_{t+1} = \phi\big(\mathcal{A}_t \cup \{\mu(\sigma(\mathcal{A}_t))\},\; V\big)
\end{equation}
The critical design choices are: the representation of candidates (code, prompts, architectures), the variation operator (LLM rewrite, crossover, mutation), the selection mechanism (tournament, roulette, novelty-based), and the archive management strategy (elitist, quality-diversity, age-based).

\textbf{Benchmark Evidence.} AlphaEvolve demonstrated the power of this loop by discovering a new matrix multiplication algorithm for $4\times 4$ complex matrices using 48 scalar multiplications---improving upon Strassen's 49-multiplication algorithm for the first time in this setting. OPRO showed that LLMs can optimize prompts for GSM8K and Big-Bench Hard, achieving improvements of 8--50\% over human-designed prompts. EvoPrompting achieved state-of-the-art on the CLRS algorithmic reasoning benchmark by evolving graph neural network architectures.

\textbf{Key Challenges.} The curse of dimensionality in the search space is the primary challenge. Even with LLM-guided mutation, the space of possible prompts, architectures, or agent designs is vast. Premature convergence---where the search settles on a local optimum---is common. DGM's approach of maintaining an open-ended archive of all candidates (rather than just the current best) is one response, but it increases computational cost. Another challenge is the cost of evaluation: each candidate must be executed in an environment, which may require significant compute time, especially for code-based agents running on complex benchmarks like SWE-Bench.

\subsubsection{The Evaluator Loop}\label{sec:loop-evaluator}

\textbf{Definition.} The Evaluator Loop provides feedback signals to the evolution process by scoring candidate artifacts against task requirements, correctness criteria, quality standards, or deployment constraints. The evaluator may be a deterministic program (unit tests, static analyzers), a model-based judge (LLM-as-a-judge), an interactive environment (web arena, code sandbox), or a human reviewer.

\textbf{Representative Systems.} FunSearch pioneered the combination of LLM generation with programmatic evaluation: LLM-generated programs are evaluated on mathematical constraints, and only those that pass both correctness and performance checks are retained in the population. Self-Debug uses execution feedback (error messages, test results) as a first-class signal for self-repair: the agent reads its own error output and generates patches. AlphaEvolve employs Gemini Pro as a deep evaluator that can assess algorithmic quality beyond simple correctness. ReVeal uses formal verification to evaluate the correctness of generated programs.

\textbf{Four-Tier Evaluation Framework.} The Evaluator Loop can be decomposed into four tiers of increasing complexity: (1) \emph{Deterministic verifiers}: unit tests, type checkers, linters, static analyzers, fuzzing, and regression tests that provide binary or scalar signals with no noise. (2) \emph{Model-based evaluators}: LLM-as-a-judge, critic agents, debate-based evaluation, and meta-judges that provide rich but noisy signals. (3) \emph{Interactive environment evaluators}: web arenas (WebArena, WebShop), code sandboxes (SWE-Bench), robot simulators, and game environments that provide grounded feedback. (4) \emph{Human evaluators}: domain experts, crowdworkers, or end users who provide preference judgments and qualitative feedback.

\textbf{Key Challenges.} The most critical challenge is \emph{evaluator degradation}. Self-Rewarding Language Models and Meta-Rewarding have demonstrated that when the evaluator is the same model (or a fine-tuned variant) as the generator, systematic biases emerge: length preference, score inflation, and reward hacking. The evaluator may also become a bottleneck: comprehensive evaluation of each candidate can be expensive, and fast approximate evaluation may miss important failure modes. Finally, the evaluator's coverage matters: a system that passes all unit tests may still fail on edge cases, security vulnerabilities, or performance regressions that the test suite does not cover.

\textbf{Connection to Safety.} The Evaluator Loop is the primary line of defense against misalignment and mis-evolution. If the evaluator can be modified by the evolving agent, or if its blind spots are exploitable, the entire evolution process becomes unreliable. Evaluator isolation---ensuring the evaluator cannot be directly modified by the agent under evaluation---is a necessary (though not sufficient) condition for safe self-evolution.

\subsubsection{The Reflection Loop}\label{sec:loop-reflection}

\textbf{Definition.} The Reflection Loop transforms execution traces, failure signals, and environmental feedback into structured knowledge that improves future behavior. Unlike the Evaluator Loop, which scores outputs, the Reflection Loop produces \emph{explanations}, \emph{principles}, \emph{insights}, or \emph{revised strategies} that modify the agent's internal state.

\textbf{Representative Systems.} Self-Refine demonstrates the simplest form: the same LLM generates an output, critiques it, and iteratively revises it within a single inference session. Reflexion extends this by persisting reflections as episodic memory: after a task failure, the agent generates a natural-language reflection describing what went wrong and how to fix it, which is retrieved on subsequent attempts. ExpeL aggregates experiences across multiple tasks to extract transferable insights---general principles that guide behavior on new tasks. STaR (Self-Taught Reasoner) uses a generate--filter--fine-tune loop: the model generates chain-of-thought reasoning, filters for correct answers, and fine-tunes on the successful traces. ACE (Adaptive Context Evolution) treats the agent's context as an evolving playbook that is incrementally refined based on task outcomes. ReasoningBank distills reasoning strategies from both successful and failed problem-solving episodes into a reusable library.

\textbf{Formal Structure.} The Reflection Loop can be formalized as a memory update operator. Let the agent's reflective memory be $m_t$ at iteration $t$. After executing a trajectory $\tau_t$ and receiving feedback $y_t$, the reflection operator $\rho$ produces an updated memory:
\begin{equation}\label{eq:reflection-loop}
m_{t+1} = \rho(m_t, \tau_t, y_t)
\end{equation}
The nature of $\rho$ varies across systems: in Reflexion, $\rho$ appends a natural-language summary to episodic memory; in ExpeL, $\rho$ extracts and generalizes insights; in STaR, $\rho$ selects correct trajectories for fine-tuning; in ACE, $\rho$ restructures the context playbook; in Memory-R1, $\rho$ is trained via RL to decide \emph{when} to write, update, delete, and retrieve memories.

\textbf{Benchmark Evidence.} Reflexion achieves 91\% pass@1 on HumanEval (compared to GPT-4's 80\%) by accumulating language-based reflections across coding episodes. STaR demonstrates that iterative self-training on filtered reasoning traces can match or exceed supervised fine-tuning on GSM8K and other math benchmarks. ExpeL shows forward transfer: insights extracted from one set of tasks improve performance on a different set. ACE shows that evolving context playbooks can improve agent performance by 15--40\% over static prompts.

\textbf{Key Challenges.} Memory quality is the central challenge. Not all reflections are equally useful; some may encode false lessons, overfit to specific task instances, or introduce contradictions with existing memory. The ``context collapse'' problem arises when too many reflections crowd the agent's context window, degrading rather than improving performance. AriadneMem addresses this with entropy-aware gating that decides which memories to retain and which to discard. Memory-R1 takes a different approach by \emph{learning} the memory management policy via reinforcement learning, rather than relying on hand-crafted heuristics.

\subsubsection{The Population Loop}\label{sec:loop-population}

\textbf{Definition.} The Population Loop maintains multiple candidate agents, strategies, or artifacts simultaneously, applying evolutionary operations---selection, mutation, recombination, and specialization---across the population. Unlike the Search Loop, which focuses on optimizing a single objective, the Population Loop emphasizes \emph{diversity}, \emph{stepping stones}, and \emph{open-ended exploration}.

\textbf{Representative Systems.} DGM (Darwin G\"{o}del Machine) maintains an ever-growing archive of agent variants: each agent can modify its own code, and the archive preserves all variants along with their lineage and performance records. AlphaEvolve uses MAP-Elites quality-diversity search to maintain a grid of high-performing solutions across multiple behavioral dimensions. Absolute Zero Reasoners employ self-play between a ``proposer'' and a ``solver,'' generating training data without any external labels. SPIRAL extends self-play to zero-sum games, discovering novel reasoning capabilities through adversarial interaction. SPIN (Self-Play Fine-Tuning) uses self-play to progressively strengthen language models: the model generates candidate responses, and a discriminator (trained on previous iterations) provides training signal.

\textbf{Key Insight: Stepping Stones.} A crucial insight from quality-diversity optimization and open-ended search is that the path to high-performing solutions often passes through candidates that are \emph{not} high-performing by the current objective. DGM's archive preserves these stepping stones, allowing the search to revisit and build upon earlier variants that may not have been the ``best'' at the time but contained useful components. This contrasts with purely elitist selection, which discards all but the current best and may converge to local optima.

\textbf{Benchmark Evidence.} DGM demonstrates that open-ended evolution with archive maintenance can produce coding agents that outperform hand-designed baselines on SWE-Bench Verified and Polyglot benchmarks. AlphaEvolve's quality-diversity search discovered new algorithms for matrix multiplication and other combinatorial problems, exceeding the capabilities of any single model generating solutions in isolation. Absolute Zero Reasoners achieve competitive performance on math and coding benchmarks using \emph{zero} external training data, relying entirely on self-play between proposal and solving agents.

\textbf{Key Challenges.} Computational cost is the primary obstacle. Maintaining and evaluating a large population of agent variants is expensive, especially when each evaluation requires running complex benchmarks. Archive management---deciding what to keep, what to discard, and how to index for efficient retrieval---is an open problem. Additionally, the population's diversity must be meaningful: maintaining many variants of the same basic approach provides little benefit.

\subsection{Cross-Loop Interactions and Composite Systems}\label{sec:theory-crossloop}

No real system implements only a single loop. The most powerful systems compose multiple loops, and the interactions between loops create emergent behaviors---both beneficial and dangerous. We identify several recurring composition patterns.

\textbf{Specification + Search.} Systems like AutoML-Agent and ADAS combine specification interpretation with search over the design space. The specification defines the task; the search explores the solution space. The interaction occurs when search results feed back into specification refinement: if no candidate satisfies the specification, the specification itself may need adjustment.

\textbf{Search + Evaluator.} This is the most fundamental composition. Every search-based system requires an evaluator to provide selection pressure. The key design decision is whether the evaluator is static (fixed benchmarks, unit tests) or co-evolving (LLM judges that are themselves being improved). AlphaEvolve uses static programmatic evaluators; Meta-Rewarding uses co-evolving model-based evaluators with a meta-judge layer to guard against evaluator degradation.

\textbf{Reflection + Population.} DGM combines reflection (agents analyze their own performance and generate self-modifications) with population-level archive maintenance. The reflection loop operates within each agent; the population loop operates across agents. The interaction creates a form of ``cultural evolution'': successful self-modifications are preserved in the archive and can be inherited by future agents.

\textbf{All Five Loops: The Self-Evolve Vision.} A complete self-evolving system would compose all five loops: specifications are parsed and refined (Loop 1), a population of candidate solutions is maintained (Loop 5), the search explores variations (Loop 2), evaluators score candidates (Loop 3), and failures trigger reflection and memory updates (Loop 4). The ``Self-Evolve'' framework proposed in this survey represents this full composition, with the added principle that each loop should be independently verifiable and auditable.

\subsection{Historical Progression and Cross-Domain Synthesis}\label{sec:theory-history}

\subsubsection{Temporal Evolution of the Field}

The field of Agent Self-Evolution has progressed through several phases, each building on the insights of its predecessors. The first phase (2022--2023) established the foundational loops: Self-Refine demonstrated iterative self-feedback at inference time; Reflexion introduced persistent language-based memory as a substitute for parameter updates; STaR showed that self-generated reasoning traces could be used for training. These early works primarily operated within the Reflection Loop, with limited integration of other loops.

The second phase (2024) saw the emergence of training-time self-improvement and multi-loop composition. Self-Rewarding Language Models combined the Evaluator Loop (model-as-judge) with the Search Loop (Iterative DPO), training the model to improve both generation and evaluation simultaneously. RISE introduced multi-turn MDP formulations for self-corrective policy learning. ExpeL and Voyager demonstrated that accumulated experience could transfer across tasks, moving beyond single-episode improvement.

The third phase (2025) brought architecture-level and code-level self-modification. ADAS automated the search over agent designs. DGM combined self-reference with population-level archive maintenance. SICA demonstrated self-modifying coding agents on real-world benchmarks. AlphaEvolve showed that LLM-guided evolutionary search could make genuine scientific discoveries. These systems compose all five loops and represent the current frontier.

The fourth phase (2025--2026) is marked by growing attention to safety, memory governance, and production readiness. Misevolution analysis identifies failure modes specific to self-evolving systems. OEP demonstrates experience-poisoning attacks. Capability Erosion studies catastrophic forgetting in self-evolving agents. AriadneMem and Memory-R1 develop sophisticated memory management policies.

\subsubsection{Cross-Domain Convergence}

A striking feature of the field is the convergence of five previously separate research traditions: (1) \textbf{AutoML + LLM} (AutoML-GPT, AutoML-Agent), where the evolved object is the ML pipeline; (2) \textbf{NAS + LLM} (EvoPrompting, LLMatic, NADER), where the evolved object is the network architecture; (3) \textbf{LLM Self-Improvement} (Self-Refine, Reflexion, STaR, SPIN), where the evolved object is the model output or weights; (4) \textbf{Evolutionary Computation + LLM} (OPRO, FunSearch, ReEvo, LLaMEA), where the evolved object is a program, prompt, or heuristic; and (5) \textbf{Agent Frameworks} (AutoGPT, MetaGPT, AutoGen, CrewAI), which provide the orchestration infrastructure.

Self-Evolving AI sits at the intersection of these five traditions. The common thread is that each tradition contributes one or more of the five loops: AutoML provides the Specification-to-Execution Loop; NAS provides the Search Loop with quality-diversity optimization; LLM self-improvement provides the Reflection Loop; EC provides the Population Loop; and all traditions contribute to the Evaluator Loop. The unified framework presented in this survey is possible precisely because these traditions are converging on the same set of design problems.

% ====================================================================
% Table 2.2: Five Evolution Loops Comparison
% ====================================================================
\begin{table}[htbp]
\centering
\small
\caption{Comparison of the five evolution loops across key dimensions. Each loop addresses a distinct aspect of self-evolution, and most real systems compose multiple loops.}
\label{tab:ch2-five-loops}
\begin{tabular}{p{2.5cm}p{2.5cm}p{2.5cm}p{2.5cm}p{3.5cm}}
\toprule
\textbf{Loop} & \textbf{Primary Target} & \textbf{Key Operation} & \textbf{Feedback Signal} & \textbf{Representative Systems} \\
\midrule
Specification-to-Execution & Task specification $\to$ runnable artifact & Interpretation, generation, refinement & Execution trace, specification satisfaction & AutoML-GPT, AutoML-Agent, ADAS \\
Search & Candidate space exploration & Variation, selection, archive management & Fitness scores, novelty, diversity & OPRO, AlphaEvolve, EvoPrompting, OpenEvolve \\
Evaluator & Quality assessment & Scoring, verification, calibration & Correctness, performance, safety compliance & FunSearch, Self-Debug, AlphaEvolve (Pro evaluator) \\
Reflection & Internal state update & Memory writing, insight extraction, principle learning & Failure explanations, success patterns & Reflexion, ExpeL, STaR, ACE, ReasoningBank \\
Population & Multi-candidate maintenance & Selection, recombination, specialization, archive growth & Multi-objective fitness, stepping stones & DGM, MAP-Elites, Absolute Zero, SPIRAL, SPIN \\
\bottomrule
\end{tabular}
\end{table}
